% Generated draft: requires mathematical and source review.
\documentclass[11pt]{article}
\usepackage[margin=25mm]{geometry}
\usepackage{fontspec}
\setmainfont{Latin Modern Roman}
\usepackage{amsmath,amssymb,mathtools,unicode-math}
\setmathfont{Latin Modern Math}
\usepackage{enumitem,needspace,xurl,hyperref}
\setlength{\parindent}{0pt}
\setlength{\parskip}{5pt}
\setlength{\emergencystretch}{3em}
\newcommand{\N}{\mathbb{N}}
\newcommand{\Z}{\mathbb{Z}}
\newcommand{\Q}{\mathbb{Q}}
\newcommand{\R}{\mathbb{R}}
\newcommand{\C}{\mathbb{C}}
\newcommand{\F}{\mathbb{F}}
\newcommand{\A}{\mathbb{A}}
\newcommand{\PP}{\mathbb{P}}
\newcommand{\E}{\mathbb{E}}
\newcommand{\eps}{\varepsilon}
\newcommand{\dd}{\,\mathrm{d}}
\begin{document}
\section*{UnsolvedMath ID: 30001966}
\textbf{Source problem code:} catalogue entry
\subsection{Definitions and mathematical statement}
The supplied record concerns mathematical instruments and their accompanying texts. In the question, \emph{explicit knowledge} refers to knowledge selected for explanation, whereas \emph{tacit knowledge} refers to knowledge left unexpressed. These are descriptive glosses; the source record supplies no precise criteria for distinguishing the two.

The source question is:
\begin{quote}
In mathematical instruments and their accompanying texts, what knowledge was selected to be explained and what knowledge remained tacit?
\end{quote}

This is a historiographical inquiry. The supplied record does not specify a corpus of instruments and texts, a historical period, a geographical scope, or criteria for identifying knowledge and determining whether it remained tacit. It therefore does not determine a mathematical proposition with specified objects, hypotheses, and quantifiers. A precise mathematical statement cannot be recovered from the supplied evidence without adding unsupported assumptions. The record explicitly classifies the statement as \texttt{unrecoverable} and describes it as a historical, pedagogical, or methodological discussion prompt.
\subsection{Short English statement}
The question asks what makers and writers chose to explain about mathematical instruments and what knowledge they left tacit. The supplied evidence supports a historical research question, but does not specify a mathematical proposition or a sufficiently defined historical corpus.
\subsection{Sources}
\begin{itemize}
\item Local UnsolvedMath registry, \texttt{data/unsolvedmath/problems.json}, permanent ID \texttt{30001966}, source code \texttt{catalogue entry}, title \emph{Explicit and Tacit Knowledge in Mathematical Instruments}. The supplied record lists 2012 as the proposed year and identifies no individual proposer.
\item \emph{Explicit Versus Tacit Knowledge in Mathematics} (2012), Oberwolfach Reports, pp.\ 131--210, DOI \texttt{10.4171/owr/2012/04}. Supplied source reference and literature-source link, labeled ``original Oberwolfach report'': \url{https://doi.org/10.4171/owr/2012/04}. The background describes Oberwolfach Reports as documenting mathematical workshops and their open-problem discussions.
\item Supplied statement-verification note: the workshop report treats this question as a historical, pedagogical, or methodological discussion prompt rather than an unresolved mathematical proposition. No separate clean statement is supplied; the original question is retained for provenance. This characterization has not been independently verified here.
\item Literature assessment dated exactly \texttt{2026-08-21}: the supplied curation labels the question ``open'' and describes it as an open historiographical question about what knowledge was made explicit, with no unique settled answer asserted by the literature. This is supplied dated curation, not an independently verified literature assessment. The background requests independent verification of time-sensitive status and cited claims.
\end{itemize}
\end{document}
